% Generated draft: requires mathematical and source review.
\documentclass[11pt]{article}
\usepackage[margin=25mm]{geometry}
\usepackage{fontspec}
\setmainfont{Latin Modern Roman}
\usepackage{amsmath,amssymb,mathtools,unicode-math}
\setmathfont{Latin Modern Math}
\usepackage{enumitem,needspace,xurl,hyperref}
\setlength{\parindent}{0pt}
\setlength{\parskip}{5pt}
\setlength{\emergencystretch}{3em}
\newcommand{\N}{\mathbb{N}}
\newcommand{\Z}{\mathbb{Z}}
\newcommand{\Q}{\mathbb{Q}}
\newcommand{\R}{\mathbb{R}}
\newcommand{\C}{\mathbb{C}}
\newcommand{\F}{\mathbb{F}}
\newcommand{\A}{\mathbb{A}}
\newcommand{\PP}{\mathbb{P}}
\newcommand{\E}{\mathbb{E}}
\newcommand{\eps}{\varepsilon}
\newcommand{\dd}{\,\mathrm{d}}
\begin{document}
\section*{UnsolvedMath ID: 30000528}
\textbf{Source problem code:} catalogue entry
\subsection{Definitions and mathematical statement}
Let \(k\) be a \(p\)-adic field, in the customary sense of a finite extension of \(\Q_p\), and let \(V/k\) be smooth and projective. Retain \(\varphi_V\) for the homomorphism in the source and put \(K_V=\ker(\varphi_V)\). The source does not state its domain or codomain, so the group to be decomposed is not fully specified by the displayed symbol alone.

The conjectural conclusion is that there exist subgroups \(F_V,D_V\subseteq K_V\) such that
\[
 K_V=F_V\oplus D_V,\qquad |F_V|<\infty,\qquad
 nD_V=D_V\quad\text{for every integer }n\geq1.
\]
Thus \(D_V\) is a divisible abelian group: every element has an \(n\)-th additive root for every positive \(n\). The assertion is not merely divisibility away from \(p\), does not require uniqueness of roots or torsion-freeness, and gives no uniform bound on \(|F_V|\). The decomposition need not be canonical. In the abelian-group interpretation it can equivalently be expressed as finiteness of \(K_V/\operatorname{Div}(K_V)\), where \(\operatorname{Div}(K_V)\) is its maximal divisible subgroup, using the usual splitting property of divisible abelian groups.

The title suggests an Albanese-kernel context. Under that additional intended identification, one would normally take \(\varphi_V\) to be the Albanese map from degree-zero zero-cycles modulo rational equivalence, \(CH_0(V)^0\), to the \(k\)-points of the Albanese variety, with the appropriate connectedness and Albanese conventions. Here \(CH_0(V)\) is the group of zero-cycles modulo principal-divisor relations on curves, and \(CH_0(V)^0\) is the kernel of the degree map. This is a conditional interpretation, not a silent definition of the absent \(\varphi_V\). Smooth projective variety also leaves geometric integrality, a chosen component or rational-point/Albanese normalization conventions unstated in this record.

The literature assessment checked 2026-08-21 describes only known cases of an unspecified conjecture and conservatively labels the incomplete scope open. The current statement does provide the finite-plus-divisible claim, but still omits the essential homomorphism definition; no named known case is supplied. The open label and editorial verification note remain dated curation metadata, not independently checked current status.
\subsection{Short English statement}
The conjecture asks whether the specified kernel over a \(p\)-adic field splits as a finite abelian group plus a group divisible by every positive integer. The map itself is undefined in the source record. Its title suggests an Albanese kernel, which is stated only conditionally rather than imposed.
\subsection{Sources}
\begin{itemize}
\item Local UnsolvedMath registry, permanent ID 30000528, source code catalogue entry, ``Albanese Kernels over \(p\)-Adic Fields,'' including the original finite-plus-divisible conjecture, cleaned question and verification metadata.
\item \emph{Algebraic K-Theory (2007)}, Oberwolfach Reports, pp.\ 1941--1990, DOI 10.4171/owr/2006/32: \url{https://doi.org/10.4171/owr/2006/32}. This is the sole supplied literature reference, labelled ``original Oberwolfach report.'' The title uses 2007, while the DOI and proposed-year metadata use 2006.
\item Literature triage checked 2026-08-21; local record updated 2026-09-16. No external source has been visited or independently checked here.
\end{itemize}
\end{document}
